% =====================================================================
% experimento_pr_aislado_costo_fixed.tex
% ---------------------------------------------------------------------
% Sección autocontenida: Experimento "pr_aislado_20260531".
% Se integra en Capitulo_Experimentos.tex via \input{}.
% =====================================================================
\providecommand{\TODO}[1]{\textcolor{rojomalo}{\textbf{[TODO: #1]}}}

% =====================================================================
\section{Experimento 10: pr\_aislado\_20260531 (aislamiento de Path Relinking)}
\label{sec:exp-pr-aislado}
% =====================================================================

\subsection{Hipótesis y motivación}

Tras aislar el selector probabilístico (Experimento~\ref{sec:exp-solo-p-inter})
y el selector binario (Experimento~\ref{sec:exp-binario-capacidad}), el tercer
approach del programa con evaluador greedy nativo aísla el efecto del
\emph{Path Relinking} (PR) como mecanismo de \textbf{intensificación} guiada
hacia la mejor solución global. La pregunta empírica es:

\begin{quote}
\textit{Bajo el costo corregido, ¿cuánto gap adicional cierra el Path
Relinking cuando se añade, como única novedad, sobre un selector ya
calibrado?}
\end{quote}

El obstáculo de partida era de \textbf{ingeniería}. El Path Relinking del
primer ciclo (\texttt{path\_relinking\_20260528}) no era aislable: estaba
implementado como un \emph{kick aumentado} que (i) solo se disparaba cuando se
disparaba el kick, (ii) usaba \texttt{sys.\_getframe} para extraer
\texttt{ctx}, $\lambda$ y la mejor solución del \emph{stack} de la
metaheurística, y (iii) parcheaba \texttt{copiar\_solucion\_labels} y
\texttt{ContadorOperadores.registrar\_mejora}. Imposible compararlo de forma
limpia contra los approaches anteriores.

\subsection{Diseño: Path Relinking limpio e independiente}

Se reescribió el Path Relinking como un módulo de \textbf{funciones puras}
(\texttt{metacarp/path\_relinking\_limpio\_20260531.py}) sin
\texttt{\_getframe} ni monkey-patches, y se expuso a las metaheurísticas
mediante un \textbf{hook explícito}: un parámetro opcional
\texttt{intensificador} (con \emph{default} \texttt{None}, por lo que el
comportamiento por defecto de las cinco MH no cambia) que se añadió a las
funciones núcleo. En el punto de estancamiento ---el mismo disparador que el
kick, controlado por \texttt{max\_iter\_sin\_mejora\_kick}--- la
metaheurística, si recibe un \texttt{intensificador}, ejecuta

\[
\texttt{sol} \leftarrow \texttt{intensificador}(\texttt{sol\_actual},\ \texttt{mejor\_global},\ \texttt{ctx},\ \lambda,\ \texttt{rng},\ \texttt{encoding},\ \texttt{md})
\]

\emph{en lugar} del kick aleatorio. Toda la información (\texttt{ctx},
$\lambda$, mejor solución) viaja por argumentos explícitos.

\begin{table}[h]
\centering
\caption{Path Relinking del primer ciclo frente al PR limpio de este experimento.}
\label{tab:pr-limpio-vs-sucio}
\small
\begin{tabular}{lll}
\toprule
Aspecto & PR primer ciclo (\texttt{20260528}) & PR limpio (\texttt{20260531}) \\
\midrule
Disparo            & Acoplado al kick            & \textbf{Hook explícito en estancamiento} \\
Acceso a \texttt{ctx}/$\lambda$ & \texttt{sys.\_getframe} (stack) & \textbf{Argumentos explícitos} \\
Mejor solución     & Patch de \texttt{copiar\_solucion} & \textbf{Argumento explícito} \\
Monkey-patches     & 3 (copia, contador, kick)   & \textbf{0} \\
Aislable           & No                          & \textbf{Sí} \\
Selector base      & Fijo (binario strict)       & \textbf{Parametrizable (p\_inter / binario)} \\
\bottomrule
\end{tabular}
\end{table}

\subsection{\texorpdfstring{Path Relinking truncado: pseudocódigo}{Path Relinking truncado: pseudocodigo}}

El Algoritmo~\ref{alg:pr-aislado} formaliza el relinking truncado. Camina
\emph{greedy} desde la solución actual hacia la guía (la mejor global)
reasignando, en cada paso, la tarea cuyo movimiento minimiza el objetivo
penalizado, y devuelve el \textbf{mejor intermedio} observado (no la solución
final). Como el mejor intermedio se inicializa con la propia
\texttt{sol\_inicio}, el resultado nunca empeora el objetivo de partida.

\begin{algorithm}[h]
\caption{Path Relinking truncado (hacia la mejor solución global).}
\label{alg:pr-aislado}
\begin{algorithmic}[1]
\Require \texttt{sol\_inicio}, \texttt{sol\_guia}; contexto \texttt{ctx};
  penalización $\lambda$.
\Ensure La mejor solución intermedia del camino \texttt{sol\_inicio}$\to$\texttt{sol\_guia}.
\State $\Delta \gets$ tareas con asignación (ruta, posición) distinta entre inicio y guía
\If{$\Delta = \varnothing$} \Return copia de \texttt{sol\_inicio} \EndIf
\State $s \gets$ \texttt{sol\_inicio}; \ $s^\star \gets s$; \
  $f^\star \gets \text{costo}(s) + \lambda\,\viol(s)$
\For{$\min(|\Delta|, 50)$ pasos}
  \State elegir $t \in \Delta$ cuyo movimiento a su posición-guía minimiza
    $\text{costo} + \lambda\,\viol$
  \State aplicar el movimiento: $s \gets s'$; \ $\Delta \gets \Delta \setminus \{t\}$
  \If{$\text{costo}(s) + \lambda\,\viol(s) < f^\star$}
    \State $s^\star \gets s$; \ $f^\star \gets \text{costo}(s) + \lambda\,\viol(s)$
  \EndIf
\EndFor
\State \Return $s^\star$
\end{algorithmic}
\end{algorithm}

\subsection{Selector parametrizable y estrategia de aislamiento}

El approach permite elegir el selector base con el parámetro
\texttt{--selector}, y lo trata como una \textbf{dimensión del grid} (se corren
ambos):

\begin{itemize}
\item \texttt{p\_inter}: selector probabilístico nativo
  (Experimento~\ref{sec:exp-solo-p-inter}), con el $p_{\text{inter}}$ ya
  calibrado por MH (SA $0.5$; TS $0.4$; RTS $0.5$; ABC $0.5$; Cuckoo $0.1$).
\item \texttt{binario}: selector binario determinista
  (Experimento~\ref{sec:exp-binario-capacidad}), instalado por monkey-patch del
  selector en el \emph{worker}.
\end{itemize}

\noindent\textbf{Aislamiento.} El efecto del Path Relinking se mide como el
$\Delta$ entre el gap de \texttt{[selector + PR]} (este experimento) y el de
\texttt{[selector]} sin PR (Experimento~\ref{sec:exp-solo-p-inter} para
\texttt{p\_inter}; Experimento~\ref{sec:exp-binario-capacidad} para
\texttt{binario}), con todos los demás componentes idénticos.

\paragraph{Nota sobre la reutilización de \emph{workers}.}
Como el selector binario se instala con un monkey-patch persistente en el
proceso, y \texttt{ProcessPoolExecutor} reutiliza los \emph{workers} entre
tareas, cada tarea fija \textbf{explícitamente} su selector al iniciar
(restaurando el selector nativo para las tareas \texttt{p\_inter}); de lo
contrario una tarea \texttt{p\_inter} podría heredar el selector binario de una
tarea previa en el mismo proceso.

\subsection{Configuración fija y dimensión del selector}

Este approach NO calibra: usa la \textbf{configuración canónica fija} por MH
(la misma que los Experimentos~8 y~9; Tabla~\ref{tab:grid-pr-aislado}) y
mantiene el \textbf{selector} como única dimensión (p\_inter vs binario). Por
tanto se corren \textbf{2 configuraciones por MH} (una por selector), cada una
con 5 repeticiones $\times$ 23 instancias, sin semilla.

\begin{table}[h]
\centering
\caption{Configuración canónica fija por MH del approach
         \texttt{pr\_aislado\_20260531} (idéntica a los Experimentos~8 y~9;
         el resto de parámetros es instance-aware).}
\label{tab:grid-pr-aislado}
\small
\begin{tabular}{lll}
\toprule
MH & Parámetro libre & Valor fijo \\
\midrule
SA            & \texttt{alpha}          & $0.90$ \\
TS simple     & \texttt{tabu\_tenure}   & $25$ \\
RTS           & \texttt{factor\_aumento}& $1.2$ \\
ABC simple    & \texttt{num\_fuentes}   & $30$ \\
Cuckoo        & \texttt{pa\_abandono}   & $0.15$ \\
\bottomrule
\end{tabular}
\end{table}

\begin{itemize}
\item Para el selector \texttt{p\_inter} se usa además el $p_{\text{inter}}$
  canónico por MH (SA $0.5$; TS $0.4$; RTS $0.5$; ABC $0.5$; Cuckoo $0.1$);
  el selector \texttt{binario} lo ignora.
\item Umbral de disparo de PR fijo en $30$ (niveles para SA, iteraciones para el
  resto); no se calibra, para no confundir el efecto de PR con su cadencia.
\end{itemize}

\subsection{Configuración y parámetros instance-aware}

Idéntica a los Experimentos~\ref{sec:exp-solo-p-inter}
y~\ref{sec:exp-binario-capacidad} salvo por la activación del Path Relinking:
operadores completos (\texttt{OPERADORES\_POPULARES}, 9), $\lambda$ por defecto
instance-aware (\texttt{lambda\_capacidad=None}), parámetros principales
instance-aware (\texttt{None}), 5 repeticiones en la fase final, sin semilla
fija, y reheat de SA acotado (\texttt{patience=10}, \texttt{reheat\_factor=0.5},
\texttt{max\_reheats\_sin\_mejora=10}). El kick aleatorio queda \emph{sustituido}
por el Path Relinking en el punto de estancamiento.

\subsection{Reproducibilidad: scripts e infraestructura}

\begin{itemize}
\item \textbf{PR limpio}: \texttt{metacarp/path\_relinking\_limpio\_20260531.py}
  (funciones puras \texttt{path\_relinking\_labels} / \texttt{path\_relinking\_ids}
  y los \emph{hooks} \texttt{hook\_pr\_labels} / \texttt{hook\_pr\_ids}).
\item \textbf{Hook} \texttt{intensificador}: parámetro opcional añadido a las 5
  metaheurísticas núcleo y a sus envoltorios \texttt{*\_desde\_instancia}
  (aditivo, \emph{default} \texttt{None}).
\item \textbf{Orquestador}: \texttt{scripts/\_pr\_aislado\_20260531\_common.py}
  (config fija, dimensión \texttt{--selector}, consolidación).
\item \textbf{Runners}: \texttt{scripts/run\_<mh>\_pr\_aislado\_20260531.py};
  \textbf{lanzador}: \texttt{scripts/run\_all\_pr\_aislado\_20260531.sh}.
\end{itemize}

La salida se escribe en
\texttt{experimentos\_costo\_fixed/<mh>\_pr\_aislado\_<YYYYMMDD-HHMM>/} con
\texttt{grid\_resumen.csv}, \texttt{mejor\_config.json} (mejor config por
selector) y \texttt{final/<mh>\_pr\_<selector>\_<instancia>.csv}.

\subsection{Resultados}

\begin{observacion}[Corrida completada]
El experimento \texttt{pr\_aislado\_20260531} se ejecutó sobre las 23 instancias
con la configuración canónica fija (la misma por MH que los Experimentos~8 y~9),
5 repeticiones por instancia, para cada uno de los dos selectores. No hay fase
de calibración: la única variable frente al Experimento~8/9 es la activación del
Path Relinking.
\end{observacion}

\begin{table}[h]
\centering
\caption{Resultados de \texttt{pr\_aislado\_20260531} (config fija, 5 reps
         $\times$ 23 instancias por MH y selector). Gap respecto al BKS con
         evaluador greedy nativo.}
\label{tab:pr-aislado-resultados}
\small
\begin{tabular}{llcccc}
\toprule
MH & Selector & Param.\ libre* & Valor* & Gap medio (\%) & $N$ inst.\ en BKS \\
\midrule
SA           & p\_inter & \texttt{alpha}           & 0.90 & \textbf{3.10}  & 18/23 \\
SA           & binario  & \texttt{alpha}           & 0.90 & 20.83 & 0/23 \\
TS simple    & p\_inter & \texttt{tabu\_tenure}    & 25   & \textbf{2.92}  & 12/23 \\
TS simple    & binario  & \texttt{tabu\_tenure}    & 25   & 17.56 & 1/23 \\
RTS          & p\_inter & \texttt{factor\_aumento} & 1.2  & \textbf{3.72}  & 14/23 \\
RTS          & binario  & \texttt{factor\_aumento} & 1.2  & 12.35 & 3/23 \\
ABC simple   & p\_inter & \texttt{num\_fuentes}    & 30   & \textbf{3.69}  & 13/23 \\
ABC simple   & binario  & \texttt{num\_fuentes}    & 30   & 5.07  & 7/23 \\
Cuckoo       & p\_inter & \texttt{pa\_abandono}    & 0.15 & \textbf{12.63} & 4/23 \\
Cuckoo       & binario  & \texttt{pa\_abandono}    & 0.15 & 18.12 & 1/23 \\
\bottomrule
\end{tabular}
\end{table}

\begin{table}[h]
\centering
\caption{Aislamiento del Path Relinking: gap medio (\%) SIN PR (Exp.~8/9)
         frente a CON PR (Exp.~10), por MH y selector. $\Delta$ negativo = PR
         mejora.}
\label{tab:pr-aislado-delta}
\small
\begin{tabular}{llrrr}
\toprule
MH & Selector & Sin PR & Con PR & $\Delta$ (pp) \\
\midrule
SA        & p\_inter & 2.89  & 3.10  & $+0.21$ \\
SA        & binario  & 20.36 & 20.83 & $+0.47$ \\
TS simple & p\_inter & 5.76  & 2.92  & $\mathbf{-2.84}$ \\
TS simple & binario  & 18.41 & 17.56 & $-0.85$ \\
RTS       & p\_inter & 4.97  & 3.72  & $-1.25$ \\
RTS       & binario  & 14.45 & 12.35 & $-2.10$ \\
ABC       & p\_inter & 9.64  & 3.69  & $\mathbf{-5.95}$ \\
ABC       & binario  & 10.22 & 5.07  & $\mathbf{-5.15}$ \\
Cuckoo    & p\_inter & 19.13 & 12.63 & $\mathbf{-6.50}$ \\
Cuckoo    & binario  & 19.57 & 18.12 & $-1.45$ \\
\bottomrule
\end{tabular}
\\[4pt]
\footnotesize
Columna ``Sin PR'': gap medio del Exp.~8 (\texttt{p\_inter}) y del Exp.~9
(\texttt{binario}). Columna ``Con PR'': este experimento (\texttt{final/}).
\end{table}

\paragraph{Interpretación de los resultados.}
El Path Relinking \textbf{mejora el gap en 8 de las 10 combinaciones}
(MH $\times$ selector); las dos excepciones son SA (ambos selectores), donde el
$\Delta$ es marginalmente positivo ($+0.21$ y $+0.47$\,pp, dentro del ruido
estocástico) porque SA ya estaba muy cerca del óptimo tras la calibración. El
efecto es muy heterogéneo y depende de cuánto margen tenía cada base:

\begin{enumerate}
\item \textbf{Ganancia enorme en las poblacionales y en TS.} Sobre la base
  \texttt{p\_inter}, ABC cae de $9.64$ a $3.69$\,\% ($-5.95$\,pp), Cuckoo de
  $19.13$ a $12.63$\,\% ($-6.50$\,pp) y TS de $5.76$ a $2.92$\,\% ($-2.84$\,pp,
  el mejor resultado del programa).
  La intensificación guiada hacia el mejor global compensa la dispersión propia
  de los métodos poblacionales y la falta de memoria de intensificación de TS.
\item \textbf{SA ya saturado; RTS mejora moderada.} SA, con su segundo knob ya
  calibrado, está tan cerca del óptimo que PR no aporta (queda en $\sim$$3$\,\%).
  RTS mejora de forma consistente ($-1.25$ en \texttt{p\_inter}, $-2.10$ en
  \texttt{binario}).
\item \textbf{PR no rescata al selector binario.} Sobre la base \texttt{binario},
  PR reduce algo el gap (p.ej.\ ABC $-5.15$, RTS $-2.10$) y en ABC lo deja
  competitivo ($5.07$\,\%), pero SA y TS siguen estancados en $\sim$$20$ y
  $\sim$$18$\,\%: el confinamiento intra-ruta del selector binario
  (Experimento~\ref{sec:exp-binario-capacidad}) es un techo que la
  intensificación por PR no puede superar por sí sola.
\end{enumerate}

\noindent\textbf{Síntesis.} La mejor base sigue siendo \texttt{p\_inter}, y
añadirle Path Relinking produce los mejores resultados del programa hasta este
punto: TS $2.92$\,\%, SA $3.10$\,\%, ABC $3.69$\,\% y RTS $3.72$\,\%
(media \texttt{p\_inter+PR} $\approx 5.2$\,\% frente a $8.48$\,\% del
Experimento~8: el Path Relinking reduce el gap medio en torno a un $40$\,\%). El beneficio de PR es real y robusto, pero se \textbf{potencia con un
buen selector}: PR sobre \texttt{p\_inter} mejora mucho más que sobre
\texttt{binario} en ABC, Cuckoo y TS. Esto confirma la jerarquía del programa:
la corrección del evaluador habilita la medición, el selector
\texttt{p\_inter} fija una base competitiva, y el Path Relinking es el primer
mecanismo de intensificación que aporta una mejora sustancial sobre esa base.

% =====================================================================
\subsection{Conclusión provisional del experimento 10}
% =====================================================================

Más allá de los gaps obtenidos, este experimento aporta:
\begin{enumerate}
\item \textbf{Un Path Relinking limpio y reutilizable.} El módulo
  \texttt{path\_relinking\_limpio\_20260531} y el hook \texttt{intensificador}
  sustituyen los \texttt{sys.\_getframe} y los tres monkey-patches del primer
  ciclo por una API explícita, aplicable a las cinco metaheurísticas.
\item \textbf{Aislamiento real del mecanismo.} Al añadir PR como única novedad
  sobre un selector ya calibrado, el $\Delta$ de gap es atribuible
  exclusivamente a la intensificación por Path Relinking.
\item \textbf{Factorización selector\,$\times$\,intensificación.} El parámetro
  \texttt{--selector} permite cuantificar si el beneficio del PR es robusto al
  selector base o si interactúa con él.
\end{enumerate}
